\documentclass{article}

\usepackage{amsmath, amsthm, amssymb, amsfonts}
\usepackage{thmtools}
\usepackage{graphicx}
\usepackage{setspace}
\usepackage{geometry}
\usepackage{float}
\usepackage{hyperref}
\usepackage[utf8]{inputenc}
\usepackage[english]{babel}
\usepackage{framed}
\usepackage[dvipsnames]{xcolor}
\usepackage{tcolorbox}
\usepackage{csquotes}
\usepackage{enumitem}
\usepackage{tikz}
\colorlet{LightGray}{White!90!Periwinkle}
\colorlet{LightOrange}{Orange!15}
\colorlet{LightGreen}{Green!15}

\newcommand{\HRule}[1]{\rule{\linewidth}{#1}}

\setstretch{1.2}
\geometry{
    textheight=9in,
    textwidth=5.5in,
    top=1in,
    headheight=12pt,
    headsep=25pt,
    footskip=30pt
}

% ------------------------------------------------------------------------------

\begin{document}

% ------------------------------------------------------------------------------
% Cover Page and ToC
% ------------------------------------------------------------------------------

\title{ \normalsize \textsc{}
		\\ [2.0cm]
		\HRule{1.5pt} \\
		\LARGE \textbf{\uppercase{Notes on Measure Theory by Cohn}
		\HRule{2.0pt} \\ [0.6cm] \LARGE{} \vspace*{10\baselineskip}}
		}
\date{}
\author{\textbf{Quinten Grubbs} \\ 
		University of Missouri-St. Louis \\
		December 15 2025}

\maketitle
\newpage

\tableofcontents
\newpage

% ------------------------------------------------------------------------------

\section{Measures}
\subsection{Algebras and Sigma-Algebras}
\textbf{Definition: } Let \(X\) be an arbitrary set. A collection \(\mathcal{A}\) of subsets of \(X\) is an \textit{algebra} on \(X\) if
\begin{enumerate}
    \item \(X\in\mathcal{A}\)
    \item for each set \(A\) that belongs to \(\mathcal{A}\), the set \(A^c\in \mathcal{A}\)
    \item for each finite sequence \(A_1,...,A_n\) of sets that belong to \(\mathcal{A}\), the set \(\cup^n_{i=1} A_i\) belongs to \(\mathcal{A}\)
    \item for each finite sequence \(A_1,...,A_n\) of sets that belong to \(\mathcal{A}\), the set \(\cap^n_{i=1}A_i\) belongs to \(\mathcal{A}\)
\end{enumerate}

We now define something similar called a \(\sigma\)-algebra if:\\
\begin{enumerate}
    \item \(X\in\mathcal{A}\)
    \item for each set \(A\) that belongs to \(\mathcal{A}\), the set \(A^c\) belongs to \(\mathcal{A}\)
    \item for each infinite sequence \(\{A_i\}\) of sets that belong to \(\mathcal{A}\), the set \(\cup_{i=1}^\infty A_i\) belongs to \(\mathcal{A}\)
    \item for each infinite sequence \(\{A_i\}\) of sets that belong to \(\mathcal{A}\), the set \(\cap_{i=1}^\infty A_i\) belongs to \(\mathcal{A}\)
\end{enumerate}

\begin{itemize}
    \item If \(X\) is a set and \(\mathcal{A}\) is a family of subsets of \(X\) that is closed under complementation, then \(X\) belongs to \(\mathcal{A}\) iff \(\emptyset\in\mathcal{A}\)
    \item If \(\mathcal{A}\) is a family of subsets of \(X\) that is nonempty, closed under complementation, and closed under the formation of finite or countable unions, then \(\mathcal{A}\) must contain \(X\).
\end{itemize}
Some examples of \(\sigma\)-algebras are as follows:
\begin{itemize}
    \item Let \(X\) be a set, then the powerset of \(X\) is a \(\sigma\)-algebra.
    \item Let \(X\) be a set, and let \(\mathcal{A}=\{\emptyset, X\}\). Then \(\mathcal{A}\) is a \(\sigma\)-algebra
    \item Let \(X\) be an infinite set, and let \(\mathcal{A}\) be the collection of all finite subsets of \(X\). Then \(\mathcal{A}\) doees not contain \(X\) and is not closed under complementation; hence it is not an algebra (or a \(\sigma\)-algebra on \(X\).
\end{itemize}

\textbf{Proposition:} Let \(X\) be a set. Then the intersection of an arbitrary nonempty collection of \(\sigma\)-algebras on \(X\) is a \(\sigma\)-algebra on \(X\).\\
\textbf{Proof:} Let \(\mathcal{C}\) be a nonempty collection of \(\sigma\)-algebras on \(X\). Now, let \(\mathcal{A}\) be the intersection of the \(\sigma\)-algebras that belong to \(\mathcal{C}\). We just need to check that \(\mathcal{A}\) contains \(X\), is closed under complementation, and is closed under the formation of countable unions. The set \(X\) belongs to \(\mathcal{A}\) since each \(\sigma\) contains \(X\). Now, let \(A\in\mathcal{A}\), each \(\sigma\)-algebra that belongs to \(\mathcal{C}\) contains \(A\) and so also contains \(\mathcal{A}^c\). Finally, we suppose \(\{A_i\}\) is a sequence of sets that belong to \(\mathcal{A}\) and hence to each \(\sigma\)-algebra in \(\mathcal{C}\). Then \(\cup_i A_i\) belongs to each \(\sigma\)-algebra in \(\mathcal{C}\) and so to \(\mathcal{A}\).\\

This proposition then implies the following result which is useful for the construction of \(\sigma\)-algebras:\\
\textbf{Corollary:} Let \(X\) be a set, and let \(\mathcal{F}\) be a family of subsets of \(X\). Then there is a smallest \(\sigma\)-algebra on \(X\) that includes \(\mathcal{F}\).\\
This smallest \(\sigma\)-algebra is, as it turns out, unique. The smallest \(\sigma\)-algebra is called the \(\sigma\)-algebra generated by \(\mathcal{F}\) and is denoted by \(\sigma(\mathcal{F})\).\\

\textbf{Proof:} Let \(\mathcal{C}\) be the collection of all \(\sigma\)-algebras on \(X\) that include \(\mathcal{F}\). Then \(\mathcal{C}\) is nonempty, since it contains the \(\sigma\)-algebra that consists of all subsets of \(X\). The intersection of the \(\sigma\)-algebras that belong to \(\mathcal{C}\) is, according to the previous proposition, a \(\sigma-algebra\); it includes \(\mathcal{F}\) and is included in every \(\sigma\)-algebra in \(\mathcal{C}\)- that is, it is included in every \(\sigma\)-algebra in \(X\) that includes \(\mathcal{F}\).\\

The \textit{Borel \(\sigma\)-algebra} on \(\mathbb{R}^d\) is the \(\sigma-algebra\) on \(\mathbb{R}^d\) generated by the collection of open subsets of \(\mathbb{R}^d\); it is denoted by \(\mathcal{B}(\mathbb{R}^d\).
 \end{document}